\documentclass[10pt,a4paper]{amsart}
\usepackage[margin=19mm]{geometry}
\usepackage{amsmath,amssymb,mathtools}
\usepackage[scr=rsfs,cal=euler]{mathalfa}
\usepackage{xfrac}
\usepackage[unicode,hidelinks,bookmarks=false]{hyperref}
\newcommand{\ie}{i.e.\ }
\newcommand{\cf}{cf.\ }
\newcommand{\resp}{respectively\ }
\newcommand{\Subsection}{\subsection}
\providecommand{\nobreakdash}{\nobreak\hbox{-}}
\setlength{\emergencystretch}{2em}
\multlinegap0pt
\usepackage{bm}
\usepackage{bbm}
\usepackage{dsfont}
\usepackage{xspace}
\usepackage{cancel}
\usepackage{mathtools}
\usepackage{amsthm}
\makeatletter
\@ifundefined{theorem}{\newtheorem{theorem}{Theorem}}{}
\@ifundefined{lemma}{\newtheorem{lemma}[theorem]{Lemma}}{}
\makeatother
\def\cH{\mathcal{H}}
\title[Power saving for the Brown-Erd\H{o}s-S\'os problem]{Power saving for the Brown-Erd\H{o}s-S\'os problem}
\author{Oliver Janzer, Abhishek Methuku, Aleksa Milojevi\'c, Benny Sudakov}
\date{Source: arXiv:2311.12765v2 (CC BY 4.0)}
\begin{document}
\maketitle

\noindent\emph{Source and licence.} Power saving for the Brown-Erd\H{o}s-S\'os problem. Version arXiv:2311.12765v2 (CC BY 4.0), distributed under the Creative Commons Attribution 4.0 International licence, as recorded in the arXiv licence metadata for this exact version and shown on the abstract page https://arxiv.org/abs/2311.12765v2. Full rights notice: licenses/arxiv-2311.12765-CC-BY-4.0.txt.\par\smallskip

\noindent\textit{Editorial scope.} Counted unit: the Lemma whose source label is \texttt{lemma:sunflower} in the article (source0.tex lines 444--446, source label \texttt{lemma:sunflower}), delivered together with the article's own complete original proof of it (source0.tex lines 447--468, one \texttt{proof} environment of 22 lines). No mathematics is written, repaired or re-derived: statement and proof are the article's own text line for line. Required context retained uncounted: none. Every definition, display and label of those lines is retained unchanged. Omitted from this package and not redistributed: 1--443, 469--641. Nothing inside a retained excerpt is omitted or altered: every retained body is delivered exactly as the article's lines are written, so no omission receipt of any kind accompanies this package. Citations: the retained excerpts carry no citation command at all, so no bibliography entry of the article is transcribed and no reference is redistributed. Reference adaptation: none was needed and none was made; every reference inside the delivered unit resolves inside it, so no unresolved reference or citation is produced. Printed slips kept literal: no printed word, symbol, hypothesis or number of the counted statement or of its proof was altered. Layout and class: the article's own class is \texttt{daj} with packages including amsmath, amssymb, amsthm, bbm, mathtools, comment, cleveref, subfigure, tikz; the wrapper is the lane's pinned \texttt{amsart} editorial wrapper at 10pt on A4, which restates the theorem-like environments the retained text needs and adds the article's own macro definitions that the retained text needs (\texttt{\textbackslash cH}), with extra wrapper packages (bm, bbm, dsfont, xspace, cancel, mathtools). The delivered numbering is the unit's own. Assets: no retained span contains a figure environment, a graphics-inclusion command, a PDF-page inclusion, a table or a TikZ picture; no image file is copied into this package, the package declares no asset dependency and no image file is redistributed with it. Licence: version arXiv:2311.12765v2 of this article is distributed under the Creative Commons Attribution 4.0 International licence, as recorded in the arXiv licence metadata for this exact version and shown on the abstract page https://arxiv.org/abs/2311.12765v2. A full rights notice is delivered at licenses/arxiv-2311.12765-CC-BY-4.0.txt. Characters: every retained excerpt is pure ASCII, so the delivered engine, XeLaTeX with its default Latin Modern text font, reports no missing character.
\begin{lemma}\label{lemma:sunflower}
Suppose that $F$ is an eligible hypergraph and suppose that a hypergraph $\cH$ contains an $(r, F)$-sunflower (for some $r$) with at least $e$ edges. If $e(F)\leq e/2$, then $\cH$ contains an $(e+\Delta(F), e)$-configuration.
\end{lemma}

\begin{proof}
By our assumption, $\cH$ contains an $(r, F)$-sunflower with at least $e$ edges. Since it is a sunflower, it has deficiency at most $\Delta(F)$. However, this sunflower may have too many edges and vertices, so it does not immediately give us an $(e+\Delta(F), e)$ configuration. Thus, the idea is to remove some vertices of degree $1$ from the copies of $F$ forming this sunflower and thus reduce the number of edges and vertices of the sunflower, without changing the deficiency in the process.

Suppose the copies of $F$ forming this sunflower correspond to embeddings $\varphi_1, \dots, \varphi_r$ and suppose that these copies are glued over a set $U\subset V(F)$ with $\Delta(U)\geq \Delta(F)$. We call the set $U$ the \textit{core} and  for each copy of $F$ forming this sunflower, we call the set $P=V(F)\backslash U$ a \textit{petal}. Further, let $e_U=e(F[U])$ denote the number of edges of $F$ completely contained in the core and let $e_P=e(F)-e(F[U])$ denote the number of edges incident to a vertex in the petal $P$. The assumption that $\Delta(U)\geq \Delta(F)$ implies that $e_P\geq |P|$, which means that every petal adds at least as many edges as vertices. Note that the sunflower constructed using only the copies $\varphi_1, \dots, \varphi_p$ would have $p|P|+|U|$ vertices and $pe_P+e_U$ edges. Let us fix $p$ to be the smallest number of petals one needs to attach to the core to get at least $e$ edges in the sunflower, i.e. $p=\lceil\frac{e-e_U}{e_P}\rceil$. Note that we may assume that $p\geq 3$ since $p=2$ implies that the sunflower has exactly $e$ edges (because we assumed $e(F) \le e/2$) and hence it represents an $( e+\Delta(F), e)$-configuration.

The total number of vertices in this sunflower is $p|P|+|U|$ and thus if $p|P|+|U|\leq e+\Delta(F)$, we are done since we have found an $(e+\Delta(F), e)$-configuration. Otherwise, $p|P|+|U|> e+\Delta(F)$, in which case our goal is to remove some vertices of degree 1 to obtain exactly $e+\Delta(F)$ vertices. Then, since the deficiency is still at most $\Delta(F)$, this gives us the desired $(e+\Delta(F), e)$-configuration. We claim that it is always sufficient to remove at most $|P|$ vertices of degree $1$ in order to bring the number of vertices to $e+\Delta(F)$. Indeed, note that although the number of vertices of the sunflower formed by $\varphi_1, \dots, \varphi_p$ is larger than $e+\Delta(F)$, it cannot be much larger. Namely, by the minimality of $p$ we have $(p-1) e_P+e_U\leq e$ and so 
\begin{align*}
    p|P|+|U|&= (p-1)|P|+ v(F)\leq (p-2)e_P + |P| + e(F)+\Delta(F)\\
    &= (p-1)e_P+e_U+|P|+\Delta(F)\leq e+\Delta(F)+|P|.
\end{align*}

Thus, it suffices to show that the sunflower consisting of $\varphi_1, \dots, \varphi_p$ has at least $|P|$ vertices of degree $1$. Let $v_U$ denote the number of vertices $v$ of $U$ such that $v$ has degree $1$ in $F$ and the edge incident to $v$ does not contain any vertices of $P$, and let $v_P$ be the number of vertices of $P$ which have degree $1$. Then the number of vertices of degree $1$ in the sunflower consisting of $\varphi_1,\dots,\varphi_p$ is equal to $v_U+pv_P$. Our goal is to show that $v_U+pv_P\geq |P|$. Suppose this is not the case and we have $v_U+pv_P<|P|$. Since $F$ is eligible, at most $e(F)/4$ vertices in the hypergraph $F$ can have degree more than $1$, so we have $v_P\geq |P|-e(F)/4$. We know that $p\geq 3$, so $v_U <  |P|-3v_P$.

Since $F$ is eligible, at least $3e(F)/4$ vertices in $F$ have degree $1$ and moreover, every edge of $F$ contains at most one vertex of degree $1$. So there are at least $3e(F)/4$ edges in $F$ incident to a vertex of degree $1$. Out of these, only $v_U$ are completely contained in $U$. Thus, the number of edges incident to $P$ is at least \[e_P\geq 3e(F)/4-v_U\geq 3e(F)/4-|P|+3v_P\geq 3e(F)/4-|P|+3(|P|-e(F)/4)=2|P|.\] 

In other words, every petal adds twice as many edges as vertices. We will show that this contradicts our assumption that $p|P| + |U| > e + \Delta(F)$. More precisely, we have 
\begin{align*}
p|P|+|U|\leq &\left(\frac{e-e_U}{e_P}+1\right)|P| +|U|\leq \frac{e}{e_P}|P|+|P|+|U|\\
\leq & \frac{e}{e_P}|P|+v(F)\leq \frac{e}{2}+e(F)+\Delta(F)\leq e+\Delta(F),
\end{align*}
since $e(F) \le e/2$. This completes the proof of the lemma.
\end{proof}

\end{document}
